\Def Дерево поиска является AVL-деревом, если для каждой вершины
высота ее правого и левого поддеревьев различаются не более чем на
единицу.

\Th Высота AVL-дерева равна $O(log\,N)$.

\Proof Рассмотрим такую величину как S(h) — минимальное число вершин в AVL-дереве высоты h. Заметим, что верно следующее соотношение

$$S(h + 2) = S(h + 1) + S(h) + 1$$

Докажем по индукции, что данную рекурренту можно выразить как $S(h) = F_{h+2} - 1$, где $F_n$ -- n-е число Фибоначчи.

Для h = 1 база верна.
Допустим, что $S(n) = F_{n+2} - 1$ для всех $n \leq h$, тогда
$S(h\,+\,1)\,=\,S(h)\,+\,S(h\,-\,1)\,+\,1\,=\,F_{h+2}\,-\,1\,+\,F_{h+1}\,-\,1\,+\,1\,=\,F_{h+2}\,+\,F_{h+1}\,-\,1\,=\,F_{h+3}\,-\,1$

Таким образом, $N \geq S(h) = F_{h+2}\,-\,1\,={\large\frac{\varphi^{h+2}-(-\varphi)^{-h-2}}{\sqrt{5}}}\,-\,1\,=\,\Theta(\varphi^h)\,\Rightarrow\,h=O(log\,N)$, где $\varphi = \frac{1 + \sqrt{5}}{2}$ \Endproof

\pagebreak